% Proved after full independent root and companion audits; exact candidate preserved.
\subsection{Rational first-level torsion and the exact index-two range}

Throughout the local assertions, \(p\ge5\), \(q=p^2\), \(N=q-1\),
and \(K\) is complete discretely valued of mixed characteristic
\((0,p)\), with algebraically closed residue field and \(\mu_p\subset K\).
Its absolute ramification \(e\) is prime to \(p\). Set \(s=e(p+1)\).
We use the actual good tame extension \(L/K\) of degree \(N\),
so that \(v_L(p)=eN=(p-1)s\) and \(N\mid s\).
A two-slope formal group has parameter \(D\) and valuations
\[
 0<D<pe,\qquad \alpha=s-D,\qquad\beta=D/p,
 \qquad \alpha>\beta.
 \tag{B-RD73.1}
\]
Its canonical high line has \(p-1\) nonzero points of valuation
\(\alpha\), and all other nonzero first-level points have valuation
\(\beta\). We always use actual integral good formal models,
actual generic torsion, and their actual inertia action.

\begin{lemma}[Rationality and the small high inverse branch]
\label{lem:B-RD73-high-branch}
\status{Proved; complete independent root and companion proof and exact-code audits passed.}
Let \(M\) be such a two-slope group. If its residual inertia action
is tame, then \(M[p]\) is rational over \(L\).
This is also true when \(M[p]\) has two distinct inertia-stable
\(\mathbf F_p\)-lines, without any semisimplicity assumption made
in advance. Conversely rationality over \(L\) implies \(p\mid D\).

Suppose \(M[p]\) is \(L\)-rational and \(D<e\).
For any nonzero high point \(Q\), the inverse equation \([p]X=Q\)
has exactly \(p\) roots of valuation
\[
 v_L(X)=\frac{C}{p},\qquad C=s-pD,
 \tag{B-RD73.2}
\]
and \(q-p\) roots of valuation \(D/p\).
The \(p\) higher-valuation roots form one coset of the high line.
Each such root has a Galois point field \(E/L\) of degree \(p\),
whose translation group is exactly that high line.
For every nonidentity \(g\in\operatorname{Gal}(E/L)\), its actual
lower break is
\[
 b(g)=eN.
 \tag{B-RD73.3}
\]
\end{lemma}
\begin{proof}
The high line is inertia stable. If residual inertia is tame,
its finite cyclic image is semisimple. Thus this stable
\(\mathbf F_p\)-line has a stable complementary line, and the
two characters take values in \(\mathbf F_p^\times\).
The tame extension of degree \(N\), a multiple of \(p-1\),
kills both characters. If two distinct stable lines are given
initially, the representation is already their direct sum,
with the same conclusion. In particular wild inertia acts
trivially: a \(p\)-group has no nontrivial character into
\(\mathbf F_p^\times\). These arguments use the algebraically
closed residue field to identify the full absolute Galois
action with inertia. They prove actual \(L\)-rationality,
not merely rationality of a semisimplification.
If all first-level points are \(L\)-rational, their low coordinate
valuation \(D/p\) is integral, so \(p\mid D\).

For the inverse equation, the exact universal coefficient
bounds used in Lemmas~\ref{lem:B-RD69-quotients}
and~\ref{lem:B-RD71-marked-low} give lower Newton vertices
\[
 (0,s-D),\qquad (p,(p-1)D),\qquad(q,0).
 \tag{B-RD73.4}
\]
When \(D<e\), the middle vertex is strictly below the chord:
this is equivalent to \((p+1)D<s\).
The initial segment gives \(p\) roots of valuation \((s-pD)/p\),
and the final segment gives \(q-p\) roots of valuation \(D/p\).
The potential degree-one point \((1,eN)\) and the remaining
coefficients lie above this polygon by those same bounds.
Here
\[
 C>D>0,\qquad p\nmid C,
 \qquad \alpha>C/p>D/p.
 \tag{B-RD73.5}
\]
The middle assertion follows from \(p\nmid s\).

Translation by any high first-level point preserves the larger
inverse valuation, because that point has valuation \(\alpha>C/p\).
Translation by a low point changes the valuation to \(D/p<C/p\).
For unequal positive parameter valuations, an integral formal
group law has valuation equal to their minimum; this follows
from its linear terms and integral terms of degree at least two.
Thus the larger inverse branch is exactly one coset of the
high line, containing \(p\) points.

The target \(Q\) and all first-level translations are rational
over \(L\). Every conjugate of a chosen larger-valuation root
remains in the same branch and is a translation of it by a high
point. In particular it belongs to \(L(X)\), so that field is
Galois of degree at most \(p\). The denominator \(p\) in
\(v_L(X)=C/p\), with \(p\nmid C\), forces ramification degree
at least \(p\). Hence \([E:L]=e(E/L)=p\), and all high-line
translations occur. Normalize \(v_E\) to be integral.
Then \(v_E(X)=C\), and every nonidentity displacement is a
high point of valuation \(p(s-D)\).
The exact prime-valuation part of Lemma~\ref{lem:B-RD60-break}
gives
\[
 b(g)=p(s-D)-C=(p-1)s=eN,
\]
as asserted. No upper/lower numbering conversion is used.
\end{proof}

\begin{proposition}[The high line is recovered from the next level]
\label{prop:B-RD73-high-recovery}
\status{Proved; complete independent root and companion proof and exact-code audits passed.}
Let \(M_1,M_2\) be two two-slope groups whose actual
\(p\)-torsion modules are \(L\)-rational. Every isomorphism of
their actual \(p^2\)-torsion identifies the high lines.
There is no upper bound on \(e\) in this assertion.
\end{proposition}
\begin{proof}
Let their parameters be \(D_1,D_2\).
They are divisible by \(p\) by the preceding lemma, so neither
equals \(e\). Assume the actual isomorphism \(\phi\) does not
identify the high lines.

Suppose first one parameter is below \(e\); after exchanging
the models and replacing \(\phi\) by its inverse if necessary,
take \(D_1<e\). Choose a larger-valuation inverse \(X_1\) of a
nonzero high point in the preceding lemma, and put \(X_2=\phi(X_1)\).
Its multiple \(pX_2\) is low in the second group. The low-target
inverse polygon gives
\[
 v_L(X_1)=\frac{s-pD_1}{p},\qquad
 v_L(X_2)=\frac{D_2}{p^3}.
 \tag{B-RD73.6}
\]
The actual torsion isomorphism identifies stabilizers, so their
point fields are the same \(E/L\), of degree \(p\).
Consequently
\[
 v_E(X_1)=s-pD_1,\qquad v_E(X_2)=D_2/q.
 \tag{B-RD73.7}
\]
If the second value is not integral, this already contradicts
\(X_2\in E\). Otherwise choose any nonidentity \(g\) in the
common degree-\(p\) group. It translates \(X_1\) by a high
point and \(X_2\) by a low point, because the first translation
group is exactly its high line and \(\phi\) sends that line
to a line distinct from the second high line.
The second displacement has valuation \(D_2\) in \(E\),
whereas the actual lower break is \(eN\). The unconditional
lower bound of Lemma~\ref{lem:B-RD60-break} therefore requires
\[
 D_2\ge D_2/q+eN,
 \qquad\text{equivalently }D_2\ge qe.
 \tag{B-RD73.8}
\]
This is impossible because \(D_2<pe<qe\).
The argument does not use distinct tame eigenvalues or a chosen
complement in a scalar residual representation.

It remains that both \(D_i>e\). Reorder so \(D_1\le D_2\).
Put \(A_1=s-D_1\); it is prime to \(p\).
For a high target in the first model the inverse polygon now
has one segment with \(q\) roots of valuation \(A_1/q\).
Choose one such root \(X_1\) and again set \(X_2=\phi(X_1)\).
Its target is low, hence \(v_L(X_2)=D_2/p^3\).
Both first-level modules are \(L\)-rational, so each point
field is Galois with translation group of order at most \(q\).
The prime-to-\(p\) numerator \(A_1\) forces the common field
to have degree \(q\); all first-level translations occur.
Its coordinate valuations are \(A_1\) and \(D_2/p\).
Choose a translation high in the first model, necessarily low
in the second. Exactly as in the degree-\(q\) comparison in
Proposition~\ref{prop:B-RD71-separated-high}, the exact break
and lower bound imply
\[
 q(D_2/p)\ge D_2/p+N A_1,
 \qquad D_2/p\ge A_1.
 \tag{B-RD73.9}
\]
But
\[
 A_1=s-D_1\ge s-D_2>D_2/p,
\]
the strict inequality being equivalent to \(D_2<pe\).
This is the final contradiction. The two cases exhaust the
possible parameters, proving the proposition.
\end{proof}

\begin{lemma}[Common-high recovery below the next cyclic threshold]
\label{lem:B-RD73-common-high}
\status{Proved; complete independent root and companion proof and exact-code audits passed.}
Assume \(e<p^3\). Two two-slope groups with isomorphic actual
\(p^3\)-torsion and corresponding high lines have the same
finite special inertia character.
\end{lemma}
\begin{proof}
Their common determinant and high-line character give
\[
 c_1\equiv c_2\pmod{p-1},\qquad
 c_1-B_1\equiv c_2-B_2\pmod N.
 \tag{B-RD73.10}
\]
Any prime-to-\(p\) parameter is settled by
Lemma~\ref{lem:B-RD69-rules}; this includes \(B_i=e\).
Assume both parameters are divisible by \(p\).
If both exceed \(e\), the actual high quotients have prime
parameters \(s-B_i\), and their inherited \(p^2\)-comparison
recovers equality. If the parameters lie on opposite sides of
\(e\), the high quotients have opposite divisibility, forbidden
by the same second-level rule.

It remains that \(B_i=pb_i<e\). Let \(M_i'=M_i/H_i\).
Their parameters are \(D_i=q b_i\), and the actual
\(p^3\)-comparison induces an actual \(p^2\)-comparison of
these quotients. By Lemma~\ref{lem:B-RD71-marked-low}, each
\(M_i'[p]\) has an intrinsic stable low line
\[
 J_i=\operatorname{im}(M_i[p]\longrightarrow M_i'[p]),
\]
distinct from its canonical high line. The rationality part of
Lemma~\ref{lem:B-RD73-high-branch} therefore shows that both
first-level modules are \(L\)-rational.
Proposition~\ref{prop:B-RD73-high-recovery} identifies their
high lines \(H_i'\), including when the residual representation
is scalar.

Take these second actual quotients. The surviving comparison
is exactly at level \(p\):
\[
 (M_i'/H_i')[p]
   =\{Z\in M_i'[p^2]:pZ\in H_i'\}/H_i'.
 \tag{B-RD73.11}
\]
Equivalently, in the original groups it is the quotient of
the corresponding subgroup of \(M_i[p^3]\) by the full
preimage of \(H_i'\), whose order is \(p^2\).
There is no unasserted second-level comparison after this step.
The second parameters are
\[
 C_i=pD_i\quad(D_i<e),\qquad
 C_i=s-D_i\quad(D_i>e).
 \tag{B-RD73.12}
\]
The equality \(D_i=e\) is impossible because \(p\mid D_i\)
and \(p\nmid e\).

If both \(D_i>e\), both second parameters are prime to \(p\).
The first-level recovery in Lemma~\ref{lem:B-RD71-first-level}
then gives \(D_1=D_2\), hence \(B_1=B_2\).
If \(D_1<e<D_2\), the first second-quotient module has two
distinct stable lines: its canonical high line and the
intrinsic low line
\[
 \operatorname{im}\bigl(M_1'[p]\longrightarrow(M_1'/H_1')[p]\bigr).
 \tag{B-RD73.13}
\]
This is another direct application of
Lemma~\ref{lem:B-RD71-marked-low}, now to the actual group
\(M_1'\) with parameter \(D_1<e\).
Its residual wild inertia is trivial. In contrast, the second
quotient has prime parameter \(s-D_2\); its low valuation
\((s-D_2)/p\) is not integral over \(L\), so its residual
division field is wild of degree \(p\). This contradicts the
actual first-level isomorphism (B-RD73.11).
The other mixed ordering is identical.

Finally, if both \(D_i<e\), then
\[
 0<b_i<e/q<p,
 \qquad b_i\in\{1,\ldots,p-1\}.
 \tag{B-RD73.14}
\]
Equation (B-RD73.10) gives \(b_1\equiv b_2\pmod{p-1}\).
Reduction modulo \(p-1\) is injective on this precise interval,
so \(b_1=b_2\). Thus every surviving branch has \(B_1=B_2\),
and the common original high character identifies \(c_1=c_2\).
This proves the assertion without a determinant square-class
normalization or a nonscalar residual assumption.
\end{proof}

\begin{proposition}[Whole formal character recovery below \(p^3\)]
\label{prop:B-RD73-formal}
\status{Proved; complete independent root and companion proof and exact-code audits passed.}
If \(e<p^3\), actual \(p^3\)-torsion determines the finite
special inertia character among whole potentially supersingular
height-two groups.
\end{proposition}
\begin{proof}
Proposition~\ref{prop:B-RD69-one-slope} handles two one-slope
groups and excludes mixed one-slope/two-slope comparisons,
without an upper restriction on \(e\).
For two two-slope groups with nontrivial residual wild inertia,
their high line is the common unique wild fixed line.
If residual inertia is tame, the first-level rationality lemma
gives \(M_i[p](L)=M_i[p]\), and
Proposition~\ref{prop:B-RD73-high-recovery} recovers their high
lines from the actual \(p^2\)-isomorphism already available.
In either case Lemma~\ref{lem:B-RD73-common-high} applies.
These are all residual possibilities.
\end{proof}

\begin{theorem}[The exact index-two range at prime-to-\(p\) ramification]
\label{thm:B-RD73}
\status{Proved; complete independent root and companion proof and exact-code audits passed.}
In the full B.4 monodromy setting of height two, with perfect
residue field, \(p\ge5\), \(\mu_p\subset K\), and \(p\nmid e\),
\[
 \boxed{\nu_{\min}(p,e,2)=2\quad\Longleftrightarrow\quad p<e<p^3.}
 \tag{B-RD73.15}
\]
Equivalently the range where the least index is two is
\[
 e=t(p-1),\qquad 2\le t\le p^2+p+1,
 \qquad p\nmid t.
 \tag{B-RD73.16}
\]
For \(p=5\) this includes every multiple of four from \(8\)
through \(124\) which is prime to five. At the next permitted
prime-to-five index \(e=128\), the least index is at least three.
\end{theorem}
\begin{proof}
For \(p<e<p^3\), pass to the completed maximal unramified base.
The cyclotomic character modulo \(p^3\) has order \(p^2\),
because \(p\nmid e\). The determinant argument of
Theorem~\ref{thm:B-RD69} identifies total dimensions \(0,1,2\)
and the finite determinant characters. Dimensions zero and two
are determined residually, with Cartier duality for the latter;
two systems filtered by potentially good height-one pieces are
also determined residually. A mixed whole-supersingular/filtered
comparison is impossible by Proposition~\ref{prop:B-RD63-mixed},
because \(e<p^3\). The two whole-group case is precisely
Proposition~\ref{prop:B-RD73-formal}. Hence actual \(p^3\)-torsion
determines the requested character in the full monodromy setting.
Theorem~\ref{thm:B-RD63-arc} at \(n=2\) gives unequal-character
systems with isomorphic actual \(p^2\)-torsion whenever \(e>p\).
Thus the least manuscript index is exactly two.

For completeness, \(\mu_p\subset K\) forces \(p-1\mid e\).
The only positive possibility with \(e\le p\) is \(e=p-1\),
whose exact index is one by Theorem~\ref{thm:B-RD61}.
There is no permitted equality \(e=p^3\), since \(p\nmid e\).
For \(e>p^3\), the same actual lower construction B63 at
\(n=3\) gives systems with equal actual \(p^3\)-torsion and
different characters, so the least index is at least three.
This proves the reverse implication as well.
Finally \((p-1)(p^2+p+1)=p^3-1\), giving the equivalent range
and the stated numerical examples. No exact index at \(e>p^3\)
is inferred.
\end{proof}

\begin{remark}[Scope of the rational-branch argument]
\label{rem:B-RD73-scope}
\status{Proved; complete independent root and companion proof and exact-code audits passed.}
The high-line recovery proposition is independent of an upper bound
on \(e\), but assumes actual \(L\)-rational first-level modules.
That hypothesis is proved separately where it is used: from tame
residual inertia, or from the intrinsic stable low line of a small
actual high quotient. Scalar residual representations cause no
exception, because the inverse-branch point fields recover the
high line without selecting a complementary eigenspace.
The common-high proof uses \(e<p^3\) only in the final integer
interval after two quotients; the full-monodromy proof also uses
the exact mixed whole/filtered threshold.
The comparison left after the second quotient is actual
\(p\)-torsion, and every subsequent argument is at that level.
No statement for absolute ramification divisible by \(p\),
no exact least index above the displayed range, and no higher-height
classification is asserted.

The retained exact supplement is
\path{code/reduction_descent/verify_rational_high_line.py}.
Its author run passed in \(9.558507562003797\) seconds at
22:30:37 UTC on 19 November 2026. For \(p=5,7,11,13\), it
checks all \(360\) permitted sectors in the displayed range,
\(275{,}760\) rational first-level polygons, the
\(22{,}932\) degree-\(p\) and \(252{,}828\) degree-\(q\)
high inverse branches, and \(22{,}932\) actual second-high-quotient
root-count records. Exact rational arithmetic reconstructs the
Newton hulls and checks the displacement inequalities and all
remaining parameter branches. Another \(48\) records above
\(p^3\) test only the unbounded high-line recovery proposition;
they do not extend the exact-index theorem. This finite supplement
does not replace the all-prime proof or establish existence of
formal groups with any enumerated parameter.

A fresh targeted search using two engines was retained with its
actual queries and read scopes. It located Hattori's finite-flat
ramification-character work and Hoshi's tame-blind extension
bound. Only the explicitly logged introduction, theorem-statement
and short-corollary passages were newly read; no unexamined main
proof is claimed as audited, no new result is imported from a
snippet, and no priority conclusion is inferred. The proof here
uses the already audited Newton, actual-quotient, lower-break and
mixed-threshold dependencies indicated above.
The independent root and companion audits passed for all six
assertions and the complete program. Their author-output-preserving
reruns took \(9.449625776003813\) and \(9.306663159004529\)
seconds, respectively. Exact prepromotion bytes, audit scopes,
the superseded weaker outline, and all executed code/output
remain retained.
\end{remark}
